\section{Methods}
\subsection{Study design and data source}
This was a secondary longitudinal forecasting analysis of the public SoccerMon dataset \citep{Midoglu2024,SoccerMonZenodo}. SoccerMon contains subjective reports collected from two teams in the Norwegian women's elite football league across the 2020 and 2021 seasons. Training-load reports were completed after training or activity sessions, whereas wellness reports were intended to be completed once per day. The original data collection used the PmSys athlete-monitoring system, was based on written participant consent for open publication, and was approved by the Norwegian Privacy Data Protection Authority (reference 296155) \citep{Midoglu2024}. The present analysis used de-identified public data and did not involve participant contact.

\subsection{Analytic unit, target, and eligibility}
The analytic unit was a player-day pair. For player $i$ on calendar day $t$, all predictors available no later than day $t$ were used to forecast readiness at day $t+1$ (Figure~\ref{fig:design}). The target was the next-calendar-day self-reported readiness-to-train score, recorded on a 1--10 scale. Because the public files provide calendar dates rather than complete within-day timestamps, the design represents a calendar-day forecasting protocol rather than a fully timestamp-verified real-time deployment simulation.

The primary complete-case cohort retained records with complete current-day wellness information, observed next-day readiness, and available 7-day load history, matching the history used by the primary model. It contained 14,582 eligible pairs from 50 pseudonymised players. The player-day eligibility flow and team/year-specific missingness are reported in Figure~\ref{fig:flow} and Supplementary Tables~S4--S5. Because primary eligibility required complete current-day wellness reports, the primary estimand applies to report-complete monitoring days rather than to all calendar days.

\begin{figure}[htbp]
\centering
\includegraphics[width=0.95\textwidth]{Figure1_study_design.pdf}
\caption{Study design and leakage-controlled forecasting protocol. Each record was defined at the player-day level. Variables reported no later than calendar day $t$ were used to forecast self-reported readiness at $t+1$. Population models were fitted only to source-team data from 2020. Sequential player-specific residual calibration used only residuals from earlier forecast dates in the target 2021 period.}
\label{fig:design}
\end{figure}

\subsection{Candidate predictors}
The primary comparator, \Pzero, used current-day readiness alone. The full monitoring model, \Pone, used current-day readiness; daily load; 7-day load; the number of training days in the preceding 7 days; current fatigue, mood, sleep duration, sleep quality, soreness, and stress; weekday and month; and recent wellness-reporting density. Daily load was treated as session-RPE aggregated to the player-day level. Load features were log-transformed where appropriate. We deliberately did not simultaneously include a large set of derived workload indicators, such as acute:chronic ratios, monotony, and strain, because these quantities are mathematically related to the same duration and RPE inputs and show substantial collinearity in SoccerMon \citep{Midoglu2024,Rago2020}. A reduced-monitoring model, \Ptwo, was examined in sensitivity analyses and included current readiness, daily load, fatigue, soreness, sleep duration, sleep quality, weekday, and month.

\begin{table}[htbp]
\centering
\caption{Forecasting model specifications.}
\label{tab:models}
\scriptsize
\begin{threeparttable}
\resizebox{\textwidth}{!}{%
\begin{tabular}{p{0.19\textwidth}p{0.47\textwidth}p{0.24\textwidth}p{0.16\textwidth}}
\toprule
Model & Population prediction model & Sequential player adaptation & Role \\
\midrule
\Pzero: adaptive autoregressive baseline & Ridge model using current-day readiness only & Mean residual across the preceding 28 eligible forecast dates for the same player; cold start in the target 2021 period & Primary comparator \\
\Pone: adaptive full-monitoring model & Ridge model using readiness; daily and 7-day load; training-day count; fatigue, mood, sleep duration, sleep quality, soreness, stress; weekday/month; reporting density & Same 28-observation residual adjustment as \Pzero & Primary intervention model \\
\Ptwo: adaptive reduced-monitoring model & Ridge model using readiness; daily load; fatigue, soreness, sleep duration, sleep quality; weekday/month & Same 28-observation residual adjustment as \Pzero & Secondary parsimony sensitivity analysis \\
\bottomrule
\end{tabular}%
}
\end{threeparttable}
\end{table}


\subsection{Population models, chronological tuning, and sequential personalisation}
For each evaluation scenario, linear/ridge population models were fitted only with the designated 2020 source data. Each comparator was tuned separately using three expanding chronological validation folds within the source-team 2020 data. Validation blocks began at approximately 55\%, 70\%, and 85\% of the source team's unique observed 2020 dates. Candidate penalties were $\alpha=0,0.0001,0.001,0.01,0.1,1,10,100,1000,10000,100000$, where $\alpha=0$ denotes the unpenalised linear-model limit.

Tuning evaluated the complete deployment pipeline. For each candidate $\alpha$, a population prediction was generated for a validation block and was then updated using the same past-only sequential residual procedure used in the 2021 evaluations. Residual history was reset at the start of every internal validation block to emulate the target-period cold-start condition. The selected penalty minimised validation-size-weighted MAE of these fully adapted validation predictions; no 2021 target data were used for tuning (Supplementary Table~S1). This design tunes the full deployed comparison rather than only the unadapted population component.

Continuous predictors were median-imputed and standardised using source-data parameters only. Let $\widehat{y}^{(m)}_{i,t+1}$ be the population prediction from model $m$ for player $i$ at time $t+1$. Both \Pzero{} and \Pone{} were then updated by the same online player-specific residual adjustment:
\begin{equation}
\widetilde{y}^{(m)}_{i,t+1} = \widehat{y}^{(m)}_{i,t+1} + \frac{1}{|\mathcal{H}_{i,t}|}\sum_{s \in \mathcal{H}_{i,t}}\left(y_{i,s+1} - \widehat{y}^{(m)}_{i,s+1}\right),
\end{equation}
where $\mathcal{H}_{i,t}$ included up to the 28 earlier eligible forecast dates for the same player in the 2021 target period. At cold start, when no earlier target-period residual was available, the adjustment was zero. This protocol reflects an online monitoring context in which prior observed readiness reports are available, but no future observations are used. It was specified as a light-touch updating strategy for operational monitoring, conceptually related to prediction-model updating in a new setting but not intended as a claim that a universal updating method is optimal \citep{Janssen2008}. The adaptation procedure was deliberately identical for \Pzero{} and \Pone{}, so the primary comparison isolates the incremental contribution of the added monitoring panel beyond the same personal-history information.

\subsection{Evaluation settings}
Four leakage-controlled settings were evaluated: (1) Team A 2020 for training and Team A 2021 for testing; (2) Team B 2020 for training and Team B 2021 for testing; (3) Team A 2020 for training and Team B 2021 for testing; and (4) Team B 2020 for training and Team A 2021 for testing. The cross-team analyses are described as a two-team transportability assessment under sequential player-specific adaptation, not as broad external validation. No random player-day split was used because the intended task was prospective forecasting with clustered repeated observations \citep{Bergmeir2018,Roberts2017,Debray2023TRIPODCluster}.

\subsection{Outcomes and statistical analysis}
The primary performance measure was mean absolute error (MAE). The primary estimand was $\Delta\mathrm{MAE}=\mathrm{MAE}(\Pone)-\mathrm{MAE}(\Pzero)$, so positive values favour the readiness-history baseline. Secondary diagnostics were root mean squared error, mean bias, $R^2$, the percentage of predictions within $\pm1$ readiness point, and the intercept and slope from linear regressions of observed on predicted readiness. We used 2,000 player-cluster bootstrap resamples to obtain 95\% confidence intervals for $\Delta\mathrm{MAE}$. These intervals quantify uncertainty across players within the evaluated target team; they do not represent sampling uncertainty across a population of teams. No smallest practically important MAE difference was prespecified; results therefore compare predictive error and uncertainty rather than formally test equivalence or clinical utility.

Sensitivity analyses varied the residual-history window (14, 28, and 56 prior forecast dates), evaluated \Ptwo, and fitted a missingness-aware full model using source-training median imputation plus wellness-specific missingness indicators in the broader cohort with observed current and next-day readiness and 7-day load history. Analyses were performed in Python using pandas, scikit-learn, NumPy, and Matplotlib. The full workflow, environment specification, and validation checks are archived in the versioned reproducibility package cited in the Code availability statement.
